\chapter{Funnel Synthesis}
\label{ch:funnel_synthesis}

\begin{gomdraftnotice}
This chapter is under active development and contains only a structural outline
with preliminary material. Full exposition, worked examples, proofs, and exercises
will be added in future editions.
\end{gomdraftnotice}

\epigraph{%
  You cannot control the exact trajectory a stone will take down a mountain.\\
  But you can build a trough to ensure it reaches the bottom.
}{}

\vspace{0.5cm}

\begin{laymansbox}{Building a Mathematical Trough}{math_trough}
Even if you know the mathematically perfect path for a golf swing, no human can execute it perfectly twice. Muscles twitch, timing fluctuates, and wind blows. If we want to guarantee success, we cannot just plan a single path; we need to build a "funnel" or a "tube" around that path. As long as the system stays anywhere inside this tube, it is mathematically guaranteed to reach the target. This chapter explains how to use powerful optimization algorithms to calculate the exact shape and size of this tube. We compute a boundary that the system is literally forbidden from crossing, ensuring that minor errors in the backswing are naturally corrected by the time the club hits the ball.
\end{laymansbox}

% ══════════════════════════════════════════════
\section{Beyond LQR: The Need for Guaranteed Invariance}
% ══════════════════════════════════════════════

In Chapter 4, we showed that applying time-varying Linear Quadratic Regulation (LQR) to the transverse dynamics creates a contracting envelope of stability. The Riccati equation naturally shapes this envelope into a ``funnel'' that narrows as the terminal constraint approaches. The idea of composing such funnels into motion planning trees was introduced by Tedrake \cite{Tedrake2009} (LQR-Trees) and extended by Majumdar and Tedrake \cite{MajumdarTedrake2017} using sum-of-squares verification.

However, LQR guarantees are inherently \emph{local}. They rely on the linear approximation of the transverse dynamics $\dot{\bm{\xi}} = \bm{A}_\perp(t)\bm{\xi} + \bm{B}_\perp(t)\delta\control$. For actual physical systems traversing highly nonlinear state spaces—like a golf club accelerating to $50$ m/s while subject to massive centrifugal and Coriolis forces—this linearization quickly breaks down. If a disturbance pushes the state too far from the reference trajectory $\traj^*(t)$, the nonlinearities will dominate, the LQR restoring forces will compute the wrong corrective action, and the clubhead will completely miss the ball.

We need a mathematically rigorous way to define the \emph{exact boundary} of our funnel. We must prove that for the fully nonlinear equations of motion, any state starting within the mouth of the funnel at time $t=0$ will remain inside the funnel for all $t \in [0, t_f]$.

\begin{definition}{Dynamic Funnel}{dynamic_funnel}
  Given a reference trajectory $\traj^*(t)$, a feedback controller $\control = \bm{k}(\state, t)$, and a scalar boundary function $V(\state, t)$, the set:
  \begin{equation}
    \mathcal{F}(t) = \{ \state \in \statespace \mid V(\state, t) \leq 1 \}
  \end{equation}
  constitutes a valid \textbf{dynamic funnel} if the set is positively invariant under the closed-loop dynamics. That is, if $\state(0) \in \mathcal{F}(0)$, then $\state(t) \in \mathcal{F}(t)$ for all $t > 0$.
\end{definition}

\begin{remark}{Funnels as Attainable Sets and Symplectic Structure}{funnels_attainable_sets}
  The dynamic funnel can be reinterpreted through the framework of \textbf{attainable sets} developed in Agrachev and Sachkov \cite{AgrachevSachkov2004}. At time $t$, the funnel boundary $V(\state, t) = 1$ defines the set of states that are on the margin of being able to reach the terminal manifold $\psi(\state(t_f)) = \bm{0}$ at time $t_f$. States strictly inside the funnel ($V < 1$) can reach the target with robustness to spare; states outside cannot reach it at all. As $t \to t_f$, the attainable set shrinks to a single point (or lower-dimensional manifold if the terminal constraint has slack). The Lyapunov function $V(\state, t)$ encodes this shrinking geometry. Furthermore, the Hamiltonian structure from the Pontryagin framework implies that the level sets $V = \text{const}$ are transverse to the true optimal trajectories in the sense of symplectic geometry. The funnel is thus not merely a Euclidean tube, but a geometric object reflecting the Hamiltonian symplectic structure of optimal motion.
\end{remark}

% ══════════════════════════════════════════════
\section{Time-Varying Lyapunov Functions for Trajectories}
% ══════════════════════════════════════════════

To verify invariance, we use the machinery of Lyapunov theory, adapted for trajectories rather than equilibria. We require our boundary function $V(\state, t)$ to act as a time-varying Lyapunov function.

\begin{theorem}[Funnel Invariance]\label{thm:funnel_invariance}
  The set $\mathcal{F}(t) = \{ \state \mid V(\state, t) \leq 1 \}$ is invariant if, on the boundary where $V(\state, t) = 1$, the function is strictly decreasing along the system trajectories:
  \begin{equation}
    \dot{V}(\state, t) = \frac{\partial V}{\partial t} + \nabla_{\state} V(\state, t) \cdot \bm{f}(\state, \bm{k}(\state, t)) < 0 \quad \text{for all } \state \text{ s.t. } V(\state, t) = 1.
  \end{equation}
\end{theorem}

If the derivative of the boundary function points inward everywhere on the boundary, trajectories may circulate inside the funnel, but they can never escape it. For a moving target problem, we want the spatial volume of $\mathcal{F}(t)$ to shrink as $t \to t_f$, terminating in a tight bounded region at impact.

\begin{center}
\begin{tikzpicture}[scale=1.2]
  % Reference trajectory
  \draw[very thick, trajgreen, decoration={markings,
      mark=at position 0.4 with {\arrow{Stealth[length=2mm]}},
      mark=at position 0.8 with {\arrow{Stealth[length=2mm]}}
    }, postaction={decorate}]
    (0, 0) .. controls (1, 2) and (2, 1.5) .. (3.5, 0.2);
  \node[trajgreen!80!black, font=\small\bfseries] at (1.5, 2.3) {$\traj^*(t)$};

  % Funnel boundaries - ellipses at different times
  \draw[thick, accentred, dotted] (0, 0) circle (0.6);
  \node[accentred, font=\scriptsize] at (-0.4, -0.8) {$t=0$};

  \draw[thick, accentred, dotted] (1.5, 1.8) ellipse (0.5 and 0.35);
  \node[accentred, font=\scriptsize] at (1.0, 2.5) {$t=t_1$};

  \draw[thick, accentred, dotted] (2.8, 0.8) ellipse (0.3 and 0.2);
  \node[accentred, font=\scriptsize] at (3.3, 1.2) {$t=t_2$};

  \draw[thick, accentred, dotted] (3.5, 0.2) circle (0.12);
  \node[accentred, font=\scriptsize] at (3.5, -0.35) {$t=t_f$};

  % Invariant tube boundary
  \draw[dashed, chapblue!60, line width=1pt] (0, 0.6) .. controls (1.2, 2.2) and (2.2, 1.85) .. (3.62, 0.32);
  \draw[dashed, chapblue!60, line width=1pt] (0, -0.6) .. controls (1.2, 1.4) and (2.2, 0.85) .. (3.38, 0.08);

  % Target manifold
  \draw[thick, accentred] (3.4, -0.15) -- (3.4, 0.55);
  \node[accentred, font=\small] at (3.85, 0.2) {$\psi = 0$};

  % Disturbance example
  \draw[-{Stealth[length=2mm]}, thick, red!70] (1.2, 1.3) -- (1.5, 0.9);
  \node[red!70, font=\small] at (0.9, 0.7) {disturbance};

  % Axes
  \draw[-{Stealth[length=2mm]}, thick, gray] (-0.5, -1.2) -- (4.2, -1.2) node[right, font=\small] {$t$};
  \draw[-{Stealth[length=2mm]}, thick, gray] (-0.5, -1.2) -- (-0.5, 2.8) node[above, font=\small] {$\state$};
\end{tikzpicture}
\end{center}

\begin{keyidea}{Funnels Define the Usable State Space}{usable_space}
  A computed funnel $\mathcal{F}(t)$ is not just an analysis tool; it is the ultimate judge of repeatability. It explicitly defines the exact set of initial conditions (e.g., the kinematic state at the top of the backswing) that are mathematically guaranteed to successfully satisfy the terminal constraint $\psi(\state(t_f)) = 0$ despite the presence of extreme nonlinearities and bounded disturbances. The time-varying geometry of the funnel---narrowing as $t \to t_f$---encodes the stability structure of the transverse dynamics along the reference trajectory.
\end{keyidea}

% ══════════════════════════════════════════════
\section{Sums-of-Squares (SOS) Programming}
% ══════════════════════════════════════════════

Proving that $\dot{V} < 0$ on the boundary $V=1$ for a general nonlinear function $\bm{f}$ is famously difficult. We have an infinite number of states on the boundary to check. We circumvent this by restricting our dynamics $\bm{f}$ algorithms to polynomials (via Taylor expansion or substitution) and leveraging \textbf{Sums-of-Squares (SOS) Programming}~\cite{Boyd1994, MajumdarTedrake2017}.

A polynomial $p(x)$ is a Sum of Squares if it can be written as $p(x) = \sum_{i} q_i(x)^2$ for some polynomials $q_i(x)$. If a polynomial is SOS, it is globally non-negative. Testing if a polynomial is SOS turns out to be equivalent to solving a Semidefinite Program (SDP)—a convex optimization problem that can be solved very efficiently.

We can reframe Theorem \ref{thm:funnel_invariance} using the S-procedure. To guarantee that $\dot{V}(\state, t) < 0$ whenever $V(\state, t) = 1$, we require:
\begin{equation}
  -\dot{V}(\state, t) - L(\state, t)(1 - V(\state, t)) \quad \text{is SOS}
\end{equation}
where $L(\state, t)$ is a strictly positive polynomial multiplier. If $V(\state, t) = 1$, the second term vanishes, and the SOS condition forces $-\dot{V} > 0$, achieving the invariance proof.

\subsection{Implementation Sketch: Funnel Synthesis}

The workflow below is the practical bridge from a nominal trajectory to a certifiable tube of allowable perturbations.

\begin{algorithm}[htbp]
\caption{Trajectory-Centered Funnel Synthesis}
\label{alg:ch7:funnel-synthesis}
\begin{algorithmic}[1]
\Require Nominal trajectory $\traj^*(t)$, feedback law $\control = \bm{k}(\state, t)$, polynomial closed-loop model, time grid $\{t_k\}$.
\Ensure Time-varying Lyapunov matrices $\bm{P}(t_k)$ defining a verified funnel.
\State Linearize or polynomially approximate the closed-loop dynamics around $\traj^*(t)$.
\State Initialize $\bm{P}(t_k)$ from the transverse Riccati solution or another stabilizing quadratic form.
\For{each node $t_k$}
  \State Define the candidate boundary $V_k(\state) = (\state - \traj^*(t_k))^\top \bm{P}(t_k) (\state - \traj^*(t_k))$.
  \State Introduce SOS multipliers enforcing $-\dot{V}_k - L_k(1 - V_k)$ to be nonnegative on the boundary.
  \State Add positive-definiteness constraints on $\bm{P}(t_k)$.
\EndFor
\While{the SDP remains feasible}
  \State Maximize the ellipsoid volume or a surrogate such as $-\log \det \bm{P}(t_k)$ across the grid.
  \State If infeasibility appears, shrink the candidate funnel, refine the polynomial model, or shorten the horizon.
\EndWhile
\State Interpolate the verified $\bm{P}(t_k)$ between nodes to obtain a continuous funnel approximation.
\Return The verified funnel family $\mathcal{F}(t) = \{\state \mid V(\state, t) \leq 1\}$
\end{algorithmic}
\end{algorithm}

% ══════════════════════════════════════════════
\section{Computing the Funnel for the Golf Swing}
% ══════════════════════════════════════════════

Let us return to the double pendulum of the golf swing. Our reference trajectory $\traj^*(t)$ from Chapter 6 gives us the nominal kinematics. Our transverse LQR from Chapter 4 gives us a feedback matrix $\bm{K}(t)$ designed to stabilize deviations. We wish to compute the largest possible funnel around this sequence that the closed-loop system is mathematically guaranteed to stay inside.

We define our candidate Lyapunov function as a time-varying quadratic form:
\begin{equation}
  V(\state, t) = (\state - \traj^*(t))^\top \bm{P}(t) (\state - \traj^*(t))
\end{equation}
where $\bm{P}(t)$ is a positive definite matrix that changes continuously along the trajectory. The funnel boundary is the time-varying ellipsoid where $V=1$. By altering $\bm{P}(t)$, we alter the shape and orientation of the funnel.

\begin{remark}{Value Functions and Hamilton-Jacobi-Bellman Theory}{hjb_value_functions}
  The Lyapunov function $V(\state, t)$ in funnel synthesis is deeply connected to the value function from dynamic programming, and more broadly to the \textbf{Hamilton-Jacobi-Bellman (HJB) equation} in optimal control. The value function $J^*(\state, t)$ represents the optimal cost to reach the terminal constraint from state $\state$ at time $t$. In the deterministic setting of Agrachev and Sachkov's framework \cite{AgrachevSachkov2004}, the value function satisfies the HJB partial differential equation:
  \begin{equation}
    \frac{\partial J^*}{\partial t} + \min_{\control} \left[ L(\state, \control) + \nabla_{\state} J^* \cdot \bm{f}(\state, \control) \right] = 0.
  \end{equation}
  The level sets of $J^*(\state, t) = c$ define the attainable set from which cost $c$ can still be incurred to reach the target. The funnel synthesized via Lyapunov methods is a conservative approximation of this optimal value function: it guarantees invariance via the $\dot{V} < 0$ condition, making it a certificate of feasibility. Computing the exact value function would require solving the HJB equation, which is generally difficult; the Lyapunov/SOS approach trades exactness for verifiability.
\end{remark}

\begin{example}{Volume Maximization via SDP}{vol_max}
  We transcribe the continuous time $t$ into discrete nodes $t_k$, just as we did in direct collocation. At each node, we set up an SDP that seeks to \emph{maximize} the volume of the ellipsoid defined by $\bm{P}(t_k)$, subject to the SOS invariance constraints connecting it to adjacent nodes.

  The optimizer actively computes the exact structural limits of the golf swing's passive dynamics. It discovers that early in the downswing, the funnel can be incredibly "wide" (forgiving), but as centrifugal forces build, the SOS constraints force $\bm{P}(t_k)$ to become highly elongated along the unactuated dimensions in order to maintain the invariance guarantee. The resulting geometry is a true numerical portrait of athletic consistency.
\end{example}

% ══════════════════════════════════════════════
\section{Summary}
% ══════════════════════════════════════════════

\begin{center}
\renewcommand{\arraystretch}{1.4}
\begin{tabular}{@{} l p{7.5cm} @{}}
  \toprule
  \textbf{Concept} & \textbf{Role in Funnel Synthesis} \\
  \midrule
  Lyapunov Function & Time-varying scalar $V(\state, t)$ that measures distance from the reference trajectory; boundary $V = 1$ defines the funnel surface. \\
  Positive Invariance & Mathematical property ensuring that $\state(0) \in \mathcal{F}(0)$ implies $\state(t) \in \mathcal{F}(t)$ for all $t \in [0, t_f]$. \\
  Sums-of-Squares (SOS) & Convex optimization technique to verify that $\dot{V} < 0$ on the funnel boundary, eliminating the need for numerical sampling. \\
  Semidefinite Program & Optimization problem solved to compute the time-varying matrix $\bm{P}(t)$ that maximizes funnel volume subject to SOS constraints. \\
  Transverse LQR & Provides initial feedback gains $\bm{K}(t)$ for stabilizing deviations orthogonal to the reference trajectory. \\
  Funnel Synthesis Workflow & Practical sequence of polynomial approximation, SOS certification, and ellipsoid-volume maximization around the nominal trajectory. \\
  \bottomrule
\end{tabular}
\end{center}

\vspace{0.3cm}

In Chapter 8, we will explore how to make these funnels robust to temporal variation by migrating out of the time domain entirely and into the phase domain.
